\documentclass{article}
\usepackage[backend=biber,natbib=true,style=alphabetic,maxbibnames=50]{biblatex}
\addbibresource{/home/nqbh/reference/bib.bib}
\usepackage[utf8]{vietnam}
\usepackage{tocloft}
\renewcommand{\cftsecleader}{\cftdotfill{\cftdotsep}}
\usepackage[colorlinks=true,linkcolor=blue,urlcolor=red,citecolor=magenta]{hyperref}
\usepackage{amsmath,amssymb,amsthm,enumitem,float,graphicx,mathtools,tikz}
\usetikzlibrary{angles,calc,intersections,matrix,patterns,quotes,shadings}
\allowdisplaybreaks
\newtheorem{assumption}{Assumption}
\newtheorem{baitoan}{Bài}
\newtheorem{cauhoi}{Câu hỏi}
\newtheorem{conjecture}{Conjecture}
\newtheorem{corollary}{Corollary}
\newtheorem{dangtoan}{Dạng toán}
\newtheorem{definition}{Definition}
\newtheorem{dinhly}{Định lý}
\newtheorem{dinhnghia}{Định nghĩa}
\newtheorem{example}{Example}
\newtheorem{ghichu}{Ghi chú}
\newtheorem{hequa}{Hệ quả}
\newtheorem{hypothesis}{Hypothesis}
\newtheorem{lemma}{Lemma}
\newtheorem{luuy}{Lưu ý}
\newtheorem{nhanxet}{Nhận xét}
\newtheorem{notation}{Notation}
\newtheorem{note}{Note}
\newtheorem{principle}{Principle}
\newtheorem{problem}{Problem}
\newtheorem{proposition}{Proposition}
\newtheorem{question}{Question}
\newtheorem{remark}{Remark}
\newtheorem{theorem}{Theorem}
\newtheorem{vidu}{Ví dụ}
\usepackage[margin=2cm]{geometry}
\def\labelitemii{$\circ$}
\DeclareRobustCommand{\divby}{%
	\mathrel{\vbox{\baselineskip.65ex\lineskiplimit0pt\hbox{.}\hbox{.}\hbox{.}}}%
}
\setlist[itemize]{leftmargin=*}
\setlist[enumerate]{leftmargin=*}

\author{Nguyễn Quản Bá Hồng}
\title{Final Exam {\it\&} Solution: Introduction to Mathematical Analysis Summer 2026\\Đề Thi Cuối Kỳ {\it\&} Lời Giải Nhập Môn Giải Tích Hè 2026}
\date{\today}

\begin{document}
\maketitle
\begin{center}
	\textbf{\textsf{Overview -- Tổng quan}}
\end{center}

\begin{enumerate}
	\item Thời gian thi: 2.5 hours.
	\item Được phép sử dụng tài liệu giấy không giới hạn.
	\item Nếu làm không được, viết định nghĩa sẽ được 2 điểm.
	\item Không làm được ý trước, vẫn có thể sử dụng kết quả các ý trước để làm ý sau của bài toán.
	\item Cấm sử dụng AIs và các thiết bị điện tử (ngoại trừ máy tính bỏ túi). Nếu phát hiện 0 điểm ngay lần đầu tiên \& thu bài ngay lập tức: không có cảnh cáo.
\end{enumerate}

\paragraph{Ký hiệu -- Notation.} $[n]\coloneqq\{1,2,\ldots,n\}$, $[n]_0\coloneqq[n]\cup\{0\} = \{0,1,\ldots,n\}$, $\mathbb{N}^*\coloneqq\{1,2,\ldots\}$, $\mathbb{R}_{\ne0}\coloneqq\mathbb{R}\backslash\{0\}$, $\mathbb{R}_{> 0}\coloneqq(0,\infty)$.

\paragraph{Yêu cầu -- Requirements.} Đề thi xoay quanh các hàm số đa thức bậc $\le3$: rất dễ, rất căn bản, \& cực kỳ sơ cấp (super easy \& very basic elementary functions). Làm 4 bài toán sau cho 4 hàm số: hàm hằng, hàm bậc nhất, hàm đa thức bậc 2 \& bậc 3:
\begin{align*}
	f_0:\mathbb{R}\to\mathbb{R},\quad f_0(x) &= a,\ a\in\mathbb{R}\\
	f_1:\mathbb{R}\to\mathbb{R},\quad f_1(x) &= ax + b,\ a\in\mathbb{R}_{\ne0},\ b\in\mathbb{R}\\
	f_2:\mathbb{R}\to\mathbb{R},\quad f_2(x) &= ax^2 + bx + c,\ a\in\mathbb{R}_{\ne0},\ b,c\in\mathbb{R}\\
	f_3:\mathbb{R}\to\mathbb{R},\quad f_3(x) &= ax^3 + bx^2 + cx + d,\ a\in\mathbb{R}_{\ne0},\ b,c,d\in\mathbb{R}.
\end{align*}
\textit{Bonus ($\star$)}: Làm bài toán với hàm số đa thức bậc $n$:
\begin{equation*}
	f_n:\mathbb{R}\to\mathbb{R},\quad f_n(x) = \sum_{i = 0}^n a_ix^i\in\mathbb{R}[x],\ a_n\in\mathbb{R}_{\ne0},a_i\in\mathbb{R},\ \forall i\in[n - 1]_0.
\end{equation*}

\begin{baitoan}[Limit of sequences -- Giới hạn của dãy số]
	Với mỗi $i\in[3]_0$:
	\begin{enumerate}[label=(\alph*)]
		\item {\rm(5 điểm)} Tính giới hạn $l$ của 2 dãy số $\{a_n\}_{n=1}^\infty,\{b_n\}_{n=1}^\infty$ lần lượt được xác định bởi $a_n\coloneqq f_i(n),\ b_n\coloneqq f_i\left(\frac{1}{n}\right)$.
		\item {\rm(10 điểm)} Với mỗi $l$ hữu hạn ở (a), chứng minh $\lim_{n\to\infty} b_n = l$ bằng ngôn ngữ $\varepsilon$.
		\item {\rm(10 điểm)} Tìm công thức để tính chỉ số tối ưu:
		\begin{equation*}
			N_\varepsilon^{\rm opt} \coloneqq \min \{ N \in \mathbb{N}^* \mid \forall n\ge N, |b_n - l| < \varepsilon \},\ \forall\varepsilon\in\mathbb{R}_{> 0}.
		\end{equation*}
	\end{enumerate}
\end{baitoan}
\textsf{Nhận xét}: Bài toán này được thiết kế với một ý đồ sư phạm rất xuất sắc để dạy cho sinh viên (đặc biệt là dân Khoa học Máy tính) phân biệt được 2 tư duy cốt lõi trong giải thuật:
\begin{enumerate}
	\item \textit{Đánh giá tiệm cận (Loose bound ở câu b)}: Dùng Bất đẳng thức tam giác để ``chặn thô'', tìm ra 1 ngưỡng $N_\varepsilon$ đủ tốt để chứng minh tính hội tụ một cách cực kỳ nhanh chóng và ít tốn chi phí tính toán $\mathcal{O}(1)$.
	\item \textit{Tối ưu hóa tuyệt đối (tight bound ở câu (c))}: Phải giải chính xác bất phương trình để tìm ra $N$ nhỏ nhất có thể $N_\varepsilon^{\rm opt}$. Đánh đổi lại, chi phí tính toán rất đắt đỏ do phải giải phương trình bậc cao.
\end{enumerate}

\begin{proof}[Giải]
	Ta xét lần lượt các dãy $a_n\coloneqq f_i(n)$ \& $b_n\coloneqq f_i(\frac{1}{n})$. Khi $n\to\infty$, ta có $\frac{1}{n}\to0$. Do các hàm đa thức đều liên tục trên $\mathbb{R}$, ta luôn có $\lim_{n\to\infty} b_n = \lim_{n\to\infty} f_i(\frac{1}{n}) = f_i(0)$. Cụ thể:
	\begin{itemize}
		\item Trường hợp $i = 0$, $f_0(x) = a$: Dãy $a_n = a$ nên $\lim_{n\to\infty} a_n = a$. Dãy $b_n = a$ nên $\lim_{n\to\infty} b_n = a$. Cả 2 dãy đều có giới hạn hữu hạn $l_0 = a$.
		\item Trường hợp $i = 1$, $f_1(x) = ax + b$ với $a\ne0$:
		\begin{align*}
			\lim_{n\to\infty} a_n &= \lim_{n\to\infty} (an + b) = \infty\cdot\text{sgn}(a),\\
			\lim_{n\to\infty} b_n &= \lim_{n\to\infty} \left(\frac{a}{n} + b\right) = b.
		\end{align*}
		Chỉ có dãy $b_n$ có giới hạn hữu hạn $l_1 = b$.
		\item Trường hợp $i = 2$, $f_2(x) = ax^2 + bx + c$, với $a\ne0$:
			\begin{align*}
			\lim_{n\to\infty} a_n &= \lim_{n\to\infty} (an^2 + bn + c) = \infty\cdot\text{sgn}(a),\\
			\lim_{n\to\infty} b_n &= \lim_{n\to\infty} \left(\frac{a}{n^2} + \frac{b}{n} + c\right) = c.
		\end{align*}
		Chỉ có dãy $b_n$ có giới hạn hữu hạn $l_2 = c$.
		\item Trường hợp $i = 3$, $f_2(x) = ax^3 + bx^2 + cx + d$, với $a\ne0$:
		\begin{align*}
			\lim_{n\to\infty} a_n &= \lim_{n\to\infty} (an^3 + bn^2 + cn + d) = \infty\cdot\text{sgn}(a),\\
			\lim_{n\to\infty} b_n &= \lim_{n\to\infty} \left(\frac{a}{n^3} + \frac{b}{n^2} + \frac{c}{n} + d\right) = d.
		\end{align*}
		Chỉ có dãy $b_n$ có giới hạn hữu hạn $l_3 = d$.
	\end{itemize}
	(b) \textit{Phương pháp}: Cho trước $\varepsilon\in\mathbb{R}_{> 0}$ tùy ý, ta cần tìm 1 số $N_\varepsilon\in\mathbb{Z}_{\ge1}$ sao cho $\forall n\ge N_\varepsilon\implies|b_n - l_i| < \varepsilon$. Vì không yêu cầu $N_\varepsilon$ phải nhỏ nhất (không cần tối ưu), ta sử dụng bất đẳng thức tam giác để ``chặn thô'' (loose bound) giúp việc tìm $N$ trở nên đơn giản.
	\begin{itemize}
		\item Với $i = 0$, $l_0 = a$: Xét $|b_n - l_0| = |a - a| = 0 < \varepsilon$. Bất đẳng thức này luôn đúng với mọi $\varepsilon\in\mathbb{R}_{> 0}$ \& $\forall n\in\mathbb{Z}_{\ge1}$. Ta chọn $N_\varepsilon = 1$.
		\item Với $i = 1$, $l_1 = b$: Xét $|b_n - l_1| = \left|\left(\frac{a}{n} + b\right) - b\right| = \frac{|a|}{n}$. Ta cần ép cho $\frac{|a|}{n} < \varepsilon\iff n > \frac{|a|}{n}$. Chọn $N_\varepsilon = \lfloor\frac{|a|}{\varepsilon}\rfloor + 1$. Khi đó, $\forall n\ge N_\varepsilon\implies n > \frac{|a|}{\varepsilon}\implies|b_n - l_1| < \varepsilon$.
		\item 
	\end{itemize}
\end{proof}

\begin{baitoan}[Limit of functions -- Giới hạn của hàm số]
	Với mỗi $i\in[3]_0$:
	\begin{enumerate}[label=(\alph*)]
		\item {\rm(5 điểm)} Tính giới hạn của hàm số $\lim_{x\to x_0} f_i(x)$ tại $x = x_0\in\mathbb{R}$.
		\item {\rm(10 điểm)} Chứng minh $\lim_{x\to x_0} f_i(x) = l$ với giới hạn $l$ tìm được ở câu (a) bằng ngôn ngữ $\varepsilon$-$\delta$.
		\item {\rm(10 điểm)} Tìm công thức để tính $\delta_\varepsilon^{\rm opt}$ tối ưu:
		\begin{equation*}
			\delta_\varepsilon^{\rm opt}(x_0) \coloneqq \sup \{ \delta \in \mathbb{R}_{>0} \mid \forall x \in \mathbb{R}, 0 < |x - x_0| < \delta \Rightarrow |f_i(x) - l| < \varepsilon \},\ \forall\varepsilon\in\mathbb{R}_{> 0}.
		\end{equation*}
	\end{enumerate}
\end{baitoan}

\begin{baitoan}[Derivative, differentiability \& numerical differentiation -- Đạo hàm, tính khả vi, \& xấp xỉ đạo hàm]
	Với mỗi $i\in[3]_0$:
	\begin{enumerate}[label=(\alph*)]
		\item {\rm(5 điểm)} Tính đạo hàm $f_i'(x_0)$.
		\item {\rm(10 điểm)} Tìm đạo hàm của hàm $f_i(x)$ tại điểm $x = x_0$ bằng định nghĩa (thông qua giới hạn).
		\item {\rm(10 điểm)} Xấp xỉ đạo hàm tại $x_0$ với bước nhảy $h \in \mathbb{R}_{>0}$ bằng $3$ công thức: Newton forward, Newton backward, \& Stirling.
	\end{enumerate}
\end{baitoan}

\begin{baitoan}[Antiderivative, integral, \& numerical integration{\tt/}quadrature -- Nguyên hàm, tích phân, \& xấp xỉ tích phân]
	Với mỗi $i\in[3]_0$ và $\alpha, \beta \in \mathbb{R}$ sao cho $\alpha < \beta$:
	\begin{enumerate}[label=(\alph*)]
		\item {\rm(5 điểm)} Tính nguyên hàm $\int f_i(x)\,{\rm d}x$.
		
		{\sf Hint.} Nguyên hàm luôn có cộng hằng số $C$ ở cuối. Tích phân thì không.
		\item {\rm(10 điểm)} Tính tích phân xác định $\int_\alpha^\beta f_i(x)\,{\rm d}x$. 
		\item {\rm(10 điểm)} Xấp xỉ tích phân $\int_\alpha^\beta f_i(x)\,{\rm d}x$ bằng trapezoidal rule:
		\begin{equation*}
			\int_\alpha^\beta f_i(x)\,{\rm d}x = \frac{\beta - \alpha}{2}\big(f_i(\alpha) + f_i(\beta)\big) - \frac{(\beta - \alpha)^3}{12}f_i''(\xi) \approx \frac{\beta - \alpha}{2}\big(f_i(\alpha) + f_i(\beta)\big)
		\end{equation*}
		với $\xi\in(\alpha,\beta)$ nào đó, \& đánh giá sai số.
	\end{enumerate}
\end{baitoan}

\begin{baitoan}[Bonus: Tổng hợp kiến thức]
    {\rm(20 điểm)} Cho 1 dãy số $\{a_n\}_{n=1}^\infty$ với số hạng được xác định bởi
    \begin{equation*}
        a_n = f(n) + g'(n) + \int_{a(n)}^{b(n)} h(x)\,{\rm d}x,\ \forall n\in\mathbb{N}^\star.
    \end{equation*}
    Tìm điều kiện để dãy số: (a) hội tụ. (b) bị chặn. (c) đơn điệu.
\end{baitoan}

%------------------------------------------------------------------------------%

\printbibliography[heading=bibintoc]

\end{document}